\documentclass[11pt, a4paper]{article}

% ---------- Codifica e lingua ----------
\usepackage[utf8]{inputenc}
\usepackage[T1]{fontenc}
\usepackage[italian]{babel}

% ---------- Matematica e font ----------
\usepackage{amsmath, amssymb, amsthm}
\usepackage{mathtools}
\usepackage{newpxtext, newpxmath}          
\usepackage[scale=0.95]{sourcesans}       
\usepackage[scale=0.95]{inconsolata}       
\usepackage{bm}
\usepackage{microtype}
\linespread{1.08}

% ---------- Pagina e colori ----------
\usepackage[left=2.3cm, right=2.3cm, top=2.3cm, bottom=2.3cm]{geometry}
\usepackage[table]{xcolor}

\definecolor{sfondo}{HTML}{15181E}
\definecolor{sfondobox}{HTML}{1F242C}
\definecolor{bordo}{HTML}{3A424E}
\definecolor{testo}{HTML}{E4E7EB}
\definecolor{tenue}{HTML}{9AA4B2}
\definecolor{accento}{HTML}{6CB6FF}
\definecolor{ambra}{HTML}{E8B15D}
\definecolor{verde}{HTML}{7CCB97}
\definecolor{rosso}{HTML}{EE8A84}
\definecolor{viola}{HTML}{C3A6F5}

\pagecolor{sfondo}
\color{testo}

% ---------- Altri pacchetti ----------
\usepackage{parskip}
\usepackage{graphicx}
\usepackage{forest}
\usetikzlibrary{mindmap, trees, arrows.meta, automata, positioning, shapes}
\usepackage{listings}
\usepackage{enumitem}
\usepackage{booktabs}
\usepackage{array}
\usepackage{titlesec}
\usepackage[most]{tcolorbox}

\forestset{
  default preamble={
    for tree={
      draw=gray,
      fill=rgb:red,0.2;green,0.2;blue,0.2,
      rounded corners,
      node options={align=center},
      font=\sffamily\small,
      edge={color=gray, line width=1pt}
    }
  }
}

% ---------- Titoli ----------
\renewcommand{\thesubsection}{\arabic{subsection}}

\titleformat{\section}
  {\sffamily\bfseries\LARGE\color{testo}}{}{0pt}{}
  [\vspace{2pt}{\color{accento}\titlerule[1.2pt]}]
\titleformat{\subsection}
  {\sffamily\bfseries\Large\color{testo}}{\textcolor{accento}{\thesubsection}}{0.6em}{}
\titleformat{\subsubsection}
  {\sffamily\bfseries\large\color{accento}}{}{0pt}{}

\titlespacing*{\section}{0pt}{0.5\baselineskip}{0.8\baselineskip}
\titlespacing*{\subsection}{0pt}{1.5\baselineskip}{0.4\baselineskip}
\titlespacing*{\subsubsection}{0pt}{1\baselineskip}{0.2\baselineskip}

% ---------- Elenchi ----------
\setlist{topsep=2pt, itemsep=2pt, parsep=0pt}
\setlist[itemize,1]{label=\textcolor{accento}{\textbullet}, leftmargin=1.3em}
\setlist[itemize,2]{label=\textcolor{tenue}{--}, leftmargin=1.2em}
\setlist[enumerate,1]{label=\textcolor{accento}{\sffamily\bfseries\arabic*}, leftmargin=1.3em}

% ---------- Box ----------
\tcbset{
  appunti/.style={
    enhanced, breakable=false,
    colback=sfondobox, coltext=testo,
    boxrule=0pt, frame hidden, arc=2pt,
    left=10pt, right=10pt, top=7pt, bottom=7pt,
    before skip=9pt, after skip=11pt,
    fonttitle=\sffamily\bfseries,
    detach title,
    before upper={\setlength{\parskip}{4pt}\tcbtitle\par},
    borderline west={2.5pt}{0pt}{#1},
    coltitle=#1
  }
}

\newtcolorbox{defbox}[1]{appunti=accento, title={#1}}
\newtcolorbox{esempio}[1][Esempio]{appunti=ambra, title={#1}}
\newtcolorbox{riepilogo}[1]{appunti=verde, title={#1}}
\newtcolorbox{proprietà}[1]{appunti=viola, title={#1}}

% ---------- Pseudocodice ----------
\lstdefinelanguage{pseudo}{
  sensitive=true,
  alsoletter={-?},
  morekeywords={function, returns, return, if, then, else, do, and, static, inputs, while, an, loop, for, each, in},
  emph={true, false},
  emphstyle=\color{ambra}
}

\lstdefinestyle{pseudostile}{
  language=pseudo,
  basicstyle=\ttfamily\small\color{testo},
  keywordstyle=\color{accento}\bfseries,
  columns=fullflexible,
  keepspaces=true,
  breaklines=true,
  breakindent=2.5em,
  mathescape=true,
  aboveskip=4pt, belowskip=0pt
}

\newtcblisting{codice}[1]{
  appunti=bordo, coltitle=tenue, title={#1},
  fonttitle=\sffamily\small\bfseries,
  listing only,
  listing options={style=pseudostile}
}

% ---------- Tabelle ----------
\arrayrulecolor{bordo}
\newcommand{\intest}[1]{\textcolor{accento}{\sffamily\bfseries #1}}
\newcommand{\vero}{\textcolor{verde}{V}}
\newcommand{\falso}{\textcolor{rosso}{F}}
\newcommand{\nota}[1]{\textcolor{tenue}{#1}}

\newcommand{\KB}{\mathrm{KB}}


\begin{document}

% ---------- Intestazione ----------
\begingroup
  \raggedright
  {\sffamily\bfseries\Huge Intelligenza artificiale\par}
  \vspace{4pt}
  {\sffamily\color{tenue} Giorgio Poluzzi \quad·\quad \today\par}
\endgroup
\vspace{1.6\baselineskip}

\section*{Logical Agents}

\subsection{Agente logico}

Un \textbf{agente logico} ha due componenti:

\begin{itemize}
  \item \textbf{Knowledge base (KB)}: insieme di frasi scritte in un linguaggio formale; costituisce il contenuto specifico del dominio;
  \item \textbf{Motore inferenziale}: algoritmi generici, indipendenti dal dominio, che ricavano nuove frasi a partire da quelle già presenti.
\end{itemize}

L'approccio è \textbf{dichiarativo}: si comunica all'agente cosa sa e gli si chiede cosa fare (\textbf{Tell \& Ask}).

Un agente si può descrivere a due livelli:
\begin{itemize}
  \item \textbf{livello di conoscenza}: cosa sa;
  \item \textbf{livello implementativo}: quali strutture dati e algoritmi utilizza.
\end{itemize}

\subsubsection*{Il ciclo dell'agente}

A ogni passo l'agente ripete lo stesso ciclo:
\begin{enumerate}
  \item \textbf{Tell}: comunica alla KB la percezione al tempo $t$;
  \item \textbf{Ask}: chiede alla KB quale azione compiere;
  \item \textbf{Tell}: registra nella KB l'azione scelta;
  \item incrementa $t$ e restituisce l'azione.
\end{enumerate}

Per funzionare, un agente deve saper:
\begin{itemize}
  \item rappresentare stati e azioni;
  \item integrare nuove percezioni;
  \item aggiornare la propria rappresentazione del mondo;
  \item dedurre proprietà nascoste del mondo;
  \item dedurre le azioni corrette da eseguire.
\end{itemize}

\subsection{Logica, modelli, entailment e inferenza}

\begin{defbox}{Logica}
  Linguaggio formale per rappresentare informazioni da cui trarre conclusioni. Ha due ingredienti:
  \begin{itemize}
    \item \textbf{sintassi}: quali frasi sono ben formate;
    \item \textbf{semantica}: il significato delle frasi, cioè quando una frase è vera in un certo modello.
  \end{itemize}
\end{defbox}

\begin{defbox}{Modello}
  Mondo formalizzato in cui si può valutare se una frase è vera oppure falsa.
  \begin{itemize}
    \item $m$ è un \textbf{modello} di $\alpha$ se $\alpha$ è vera in $m$;
    \item $M(\alpha)$ è l'insieme di tutti i modelli di $\alpha$.
  \end{itemize}
\end{defbox}

\begin{defbox}{Entailment (conseguenza logica)}
  $\KB \models \alpha$ (``KB implica logicamente $\alpha$'') significa che $\alpha$ è vera in \textbf{ogni} modello in cui KB è vera. In termini di insiemi:
  \[
    \KB \models \alpha \iff M(\KB) \subseteq M(\alpha)
  \]
\end{defbox}

\begin{esempio}
  Se $\KB = \{P \land Q\}$ e $\alpha = P$, allora $\KB \models \alpha$: in ogni modello in cui $P \land Q$ è vera, anche $P$ è vera.
\end{esempio}

\begin{defbox}{Inferenza}
  $\KB \vdash_i \alpha$ significa che $\alpha$ può essere derivata dalla KB usando la procedura $i$.
\end{defbox}

\textbf{Perché serve.} Le conseguenze di una KB sono tantissime e $\alpha$ è quasi impossibile da indovinare: l'\textbf{entailment} stabilisce che $\alpha$ segue dalla KB, l'\textbf{inferenza} fornisce il metodo per trovarla.

Una procedura $i$ si valuta con due proprietà:

\begin{center}
  \renewcommand{\arraystretch}{1.35}
  \begin{tabular}{>{\bfseries}l l l}
    \toprule
    \intest{Proprietà} & \intest{Significato} & \intest{Formula} \\
    \midrule
    Correttezza & deriva \textbf{soltanto} cose vere & se $\KB \vdash_i \alpha$, allora $\KB \models \alpha$ \\
    Completezza & deriva \textbf{tutto} ciò che è vero & se $\KB \models \alpha$, allora $\KB \vdash_i \alpha$ \\
    \bottomrule
  \end{tabular}
\end{center}

\clearpage
\subsection{Logica proposizionale}

È la logica più semplice.

\subsubsection*{Sintassi}

I simboli $P_1, P_2, \dots$ sono \textbf{frasi atomiche}. Date le frasi $S, S_1, S_2$, si costruiscono nuove frasi così:

\begin{center}
  \renewcommand{\arraystretch}{1.25}
  \begin{tabular}{l l}
    \toprule
    \intest{Frase} & \intest{Nome} \\
    \midrule
    $\lnot S$ & negazione \\
    $S_1 \land S_2$ & congiunzione (``e'') \\
    $S_1 \lor S_2$ & disgiunzione (``o'') \\
    $S_1 \Rightarrow S_2$ & implicazione \\
    $S_1 \Leftrightarrow S_2$ & bicondizionale (``se e solo se'') \\
    \bottomrule
  \end{tabular}
\end{center}

\subsubsection*{Semantica}

Un \textbf{modello} assegna vero (\vero) o falso (\falso) a ogni simbolo, per esempio $P_{1,2}=F$, $P_{2,2}=F$, $P_{3,1}=V$. Con 3 simboli ci sono $2^3 = 8$ modelli, enumerabili automaticamente.

\textbf{Regole di valutazione:}
\begin{itemize}
  \item $\lnot S$ è vera se $S$ è falsa;
  \item $S_1 \land S_2$ è vera se entrambe sono vere;
  \item $S_1 \lor S_2$ è vera se almeno una è vera;
  \item $S_1 \Rightarrow S_2$ è falsa \textbf{solo} quando $S_1$ è vera e $S_2$ è falsa;
  \item $S_1 \Leftrightarrow S_2$ è vera quando valgono entrambe le implicazioni ($S_1 \Rightarrow S_2$ e $S_2 \Rightarrow S_1$).
\end{itemize}

\textbf{Tavola di verità:}

\begin{center}
  \renewcommand{\arraystretch}{1.25}
  \begin{tabular}{*{2}{>{\centering\arraybackslash}p{0.9cm}}@{\hspace{0.8cm}}*{5}{>{\centering\arraybackslash}p{1.7cm}}}
    \toprule
    $S_1$ & $S_2$ & $\lnot S_1$ & $S_1 \land S_2$ & $S_1 \lor S_2$ & $S_1 \Rightarrow S_2$ & $S_1 \Leftrightarrow S_2$ \\
    \midrule
    \falso & \falso & \vero & \falso & \falso & \vero & \vero \\
    \falso & \vero & \vero & \falso & \vero & \vero & \falso \\
    \vero & \falso & \falso & \falso & \vero & \falso & \falso \\
    \vero & \vero & \falso & \vero & \vero & \vero & \vero \\
    \bottomrule
  \end{tabular}
\end{center}

\begin{esempio}[Esempio di valutazione ricorsiva]
  Con il modello di sopra:
  \[
    \lnot P_{1,2} \land (P_{2,2} \lor P_{3,1}) = V \land (F \lor V) = V \land V = V
  \]
\end{esempio}

\clearpage
\subsection{Inferenza per enumerazione}

Si enumerano i modelli con una \textbf{ricerca in profondità} (depth-first search): a ogni passo si fissa un simbolo a vero e poi a falso. $\KB \models \alpha$ vale se $\alpha$ è vera in tutti i modelli in cui è vera la KB.

\begin{codice}{Inference by enumeration 1}
function TT-Entails?(KB, $\alpha$) returns true or false
  inputs: KB, $\text{knowledge base, a sentence in prop logic}$
          $\alpha$, $\text{a query, sentence in prop logic}$

  symbols $\leftarrow$ $\text{a list of the proposition symbols in KB and }\alpha$
  return TT-Check-All(KB, $\alpha$, symbols, [])
\end{codice}

\begin{codice}{Inference by enumeration 2}
function TT-Check-All(KB, $\alpha$, symbols, model) returns true or false
  if Empty?(symbols) then
    if PL-True?(KB, model) then return PL-True?($\alpha$, model)
    else return true
  else do
    P $\leftarrow$ First(symbols); rest $\leftarrow$ Rest(symbols)
    return TT-Check-All(KB, $\alpha$, rest, Extend(P, true, model))
       and TT-Check-All(KB, $\alpha$, rest, Extend(P, false, model))
\end{codice}

\subsection{Metodi di dimostrazione}

\begin{tcbraster}[raster columns=2, raster equal height, raster column skip=10pt, raster before skip=8pt, raster after skip=10pt]
  \begin{defbox}{Model checking}
    \raggedright
    \begin{itemize}
      \item enumerazione con tavole di verità
      \item backtracking (DPLL)
      \item ricerca locale euristica, min-conflicts/hill-climbing
    \end{itemize}
  \end{defbox}
  \begin{defbox}{Applicazione delle regole di inferenza}
    \raggedright
    \begin{itemize}
      \item si generano frasi nuove sulla base di quelle vecchie
      \item una dimostrazione è una sequenza di applicazioni di regole
      \item viene richiesto di portare le frasi in forma normale
    \end{itemize}
  \end{defbox}
\end{tcbraster}

\subsection{Equivalenza logica}

\begin{defbox}{Definizione}
  Due frasi sono equivalenti se e solo se sono vere negli stessi modelli, ovvero se una implica l'altra:
  \[
    \alpha \equiv \beta \iff \alpha \models \beta \ \text{ e } \ \beta \models \alpha
  \]
\end{defbox}

\begin{center}
  \renewcommand{\arraystretch}{1.3}
  \begin{tabular}{>{$}l<{$} >{$}c<{$} >{$}l<{$} l}
    \toprule
    \multicolumn{3}{l}{\intest{Equivalenza}} & \intest{Nome} \\
    \midrule
    (\alpha \land \beta) & \equiv & (\beta \land \alpha) & commutatività di $\land$ \\
    (\alpha \lor \beta) & \equiv & (\beta \lor \alpha) & commutatività di $\lor$ \\
    \addlinespace[4pt]
    ((\alpha \land \beta) \land \gamma) & \equiv & (\alpha \land (\beta \land \gamma)) & associatività di $\land$ \\
    ((\alpha \lor \beta) \lor \gamma) & \equiv & (\alpha \lor (\beta \lor \gamma)) & associatività di $\lor$ \\
    \addlinespace[4pt]
    \lnot(\lnot\alpha) & \equiv & \alpha & eliminazione della doppia negazione \\
    \addlinespace[4pt]
    (\alpha \Rightarrow \beta) & \equiv & (\lnot\beta \Rightarrow \lnot\alpha) & contrapposizione \\
    (\alpha \Rightarrow \beta) & \equiv & (\lnot\alpha \lor \beta) & eliminazione dell'implicazione \\
    (\alpha \Leftrightarrow \beta) & \equiv & ((\alpha \Rightarrow \beta) \land (\beta \Rightarrow \alpha)) & eliminazione del bicondizionale \\
    \addlinespace[4pt]
    \lnot(\alpha \land \beta) & \equiv & (\lnot\alpha \lor \lnot\beta) & De Morgan \\
    \lnot(\alpha \lor \beta) & \equiv & (\lnot\alpha \land \lnot\beta) & De Morgan \\
    \addlinespace[4pt]
    (\alpha \land (\beta \lor \gamma)) & \equiv & ((\alpha \land \beta) \lor (\alpha \land \gamma)) & distributività di $\land$ su $\lor$ \\
    (\alpha \lor (\beta \land \gamma)) & \equiv & ((\alpha \lor \beta) \land (\alpha \lor \gamma)) & distributività di $\lor$ su $\land$ \\
    \bottomrule
  \end{tabular}
\end{center}

\subsection{Validità \& soddisfacibilità}

\begin{tcbraster}[raster columns=3, raster equal height, raster column skip=10pt, raster before skip=8pt, raster after skip=10pt]
  \begin{defbox}{Validità}
    \raggedright
    Una proposizione è valida (tautologia) se è vera in tutti i modelli.
    \begin{esempio}
      $A \lor \lnot A$ \\
      $A \Rightarrow A$ \\
      $(A \land (A \Rightarrow B)) \Rightarrow B$
    \end{esempio}
  \end{defbox}
  \begin{defbox}{Soddisfacibilità}
    \raggedright
    Una proposizione è soddisfacibile se è vera in almeno un modello.
    \begin{esempio}
      $A \lor B$ \\
      $C$
    \end{esempio}
  \end{defbox}
  \begin{defbox}{Insoddisfacibile}
    \raggedright
    Una proposizione è insoddisfacibile se non è vera in nessun modello.
    \begin{esempio}
      $A \land \lnot A$
    \end{esempio}
  \end{defbox}
\end{tcbraster}

\subsubsection*{Approccio con l'inferenza}

  Teorema di deduzione:

  \begin{itemize}
    \item   $\KB \models \alpha \iff (\KB \Rightarrow \alpha)$ è valida.
  \end{itemize}
  Refutazione:
  \begin{itemize}
    \item $\KB \models \alpha \iff (\KB \land \lnot\alpha)$ è insoddisfacibile.
  \end{itemize}

\begin{defbox}{Sistema di inferenza}
  Un sistema di inferenza è un insieme di regole $\mathcal{R}$ della forma:
  \[
    \frac{A_1 \ \dots \ A_k}{A}
    \qquad \text{con} \qquad
    \frac{\text{premesse}}{\text{conclusione}}
  \]
\end{defbox}

Si ha una \textbf{derivazione} di $A$ da un insieme di formule $\Gamma$ con un sistema inferenziale $\mathcal{R}$ ($\Gamma \vdash_{\mathcal{R}} A$) se esiste una successione $A_1, \dots, A_n$ tale che:
\begin{itemize}
  \item $A$ è $A_n$;
  \item per ogni $i \in \{1, \dots, n\}$ vale una delle due condizioni:
  \begin{enumerate}
    \item $A_i \in \Gamma$, oppure
    \item $A_i$ è una derivazione diretta di formule che la precedono all'interno della stessa successione.
  \end{enumerate}
\end{itemize}

La successione $A_1, \dots, A_n$ è una \textbf{dimostrazione} di $A$, e l'insieme di formule $\Gamma$ viene definito come insieme delle premesse (o assunzioni/ipotesi) per $A$.

\begin{proprietà}{Proprietà dei sistemi di inferenza}
  \begin{itemize}
    \item Correttezza: le conclusioni sono sempre conseguenze logiche delle premesse $\psi_1,\dots,\psi_n\models\psi$.
    \item Completezza: se $\text{H}\models\psi$, esiste allora una derivazione di $\psi$ da H.
    \item Completezza refutazionale: se $\text{H}\cup\{\lnot\psi\}$ è insoddisfacibile, ne si riesce a derivare la contraddizione.
  \end{itemize}
\end{proprietà}

\subsubsection*{Problema di deduzione}

  Si vuole dimostrare $\Gamma\models\text{A}$. Si può quindi procedere per refutazione (bisogna dimostrare che $\Gamma\lnot\text{A}$ è insoddisfacibile), o per Forward/Backward chaining, ovvero algoritmi polinomiali (funzionano solo su una classe ristretta di formule, ovvero le clausole di Horn).

  CNF - Forma Normale Congiuntiva, ovvero una congiunzione ($\wedge$) di clausole, in cui ogni clausole è una disgiunzione ($\vee$) di litterali.

  Ogni formula è convetibile in CNF.

  \subsubsection*{Risoluzione}

  Se una clausola contiene un litterale $\ell$ e un'altra contiene il suo opposto $\lnot\ell$ li elimino entrambi e unisco il resto in una nuova clausola.

  \begin{codice}{Algoritmo di risoluzione}
    function PL-Resolution(KB,$\alpha$) returns true or false
      inputs: KB,$\alpha$
      clauses $\leftarrow$ set of clauses in CNF form of KB$\wedge\lnot\alpha$
      new $\leftarrow\{\}$
      loop do
        for each C$_i$, C$_j$ in clauses do
          resolvents $\leftarrow$ PL-Resolve(C$_i$,C$_j$)
          if resolvents contains the empty clause then return true
          new $\leftarrow$ new $\cup$ resolvents
        if new $\subseteq$ clauses then return false
        clauses $\leftarrow$ clauses $\cup$ new
  \end{codice}

  \subsubsection*{Forward \& Backward chaining}
  KB = è la congiunzione delle clausole di Horn.

  Una clausola di Horn può essere di due tipi:
  \begin{itemize}
    \item un simbolo singolo, ovvero un fatto (A);
    \item una congiunzione di simboli che implica un simbolo $(\text{C}\wedge\text{D}\Rightarrow\text{B})$
  \end{itemize}
  Su queste clausole si utilizza la regola del \textbf{Modus Ponens} (completa con le clausole di Horn, e gli algoritmi che ne derivano impiegano un tempo di esecuzione lineare): se esistono $\alpha_1,\dots,\alpha_n$ e $\alpha_1\wedge\dots\wedge\alpha_n\Rightarrow\beta$, allora è possibile dedurre $\beta$.

  Si vuole quindi attivare ogni regola le cui premesse sono già note, aggiungere le conclusioni alla coda della KB e ripetere fino alla risoluzione della query.

  \begin{codice}{Forward Chanining Algorithm}
    function PL-FC-Entails? (KB,q) returns true or false
      inputs: KB, q
      local variables: count, a table indexed by clause
                       inferred, a table indexed by symbol
                       agenda, a list of symbols
      while agenda is not empty do
        p $\leftarrow$ Pop(agenda)
        unless inferred[p] do
          inferred[p] $\leftarrow$ true
          for each Horn Clause c in whose premise p appears do
            decrement count[c]
            if count[c] = 0 then do
              if Head[c] = q then return true
              Push(Head[c], agenda)
      return false
  \end{codice}

  \subsubsection*{Completezza del Forward Chaining}

  La Forward Chaining deriva ogni fatto (simbolo singolo) atomico implicato dalla Keyword Base (KB).

  L'algoritmo si ferma in un punto fisso, in cui non deduce nulla di nuovo e considera lo stato finale come un modello $m$, ovvero vero per ciò che è stato dedotto e falso per tutto il resto.

  Ogni clausola della KB è vera in $m$, infatti se una regola in una successione $\text{a}_1\dots\text{a}_\text{k}\Rightarrow\text{b}$ fosse falsa, le premesse sarebbero tutte vere e b falso, ma così l'algoritmo avrebbe ancora qualcosa da dedurre e non ci sarebbe un punto di arrivo fisso (contraddizione).

  $\Rightarrow$ quindi $m$ è un modello della KB, e se $\text{KB}\models\text{a}$, allora a è vero in ogni modello della KB, compreso $m$.

  \subsubsection*{Backward Chaining}

  L'idea è quella di lavorare all'indietro partendo dalla query q, altrimenti (qualora non funzioni o diventi troppo complesso) si cerca di trovare una regola che possa concludere q e provi a dimostrare ricorsivamente tutte le sue premesse.

  È una buona pratica evitare i cicli (controllare sempre se un subgoal è già nello stack di goal) e il lavoro ripetuto (ricordare quali subgoal sono stati approvati e dimostrati e quali sono invece falliti).

  \subsubsection*{Confronto tra backward e forward chaining}

  Il forward chaining è manovrato dai dati, richiede spesso tanto lavoro inutile perchè deduce tutto, anche se inutile per la query.

  Il backward chaining invece è più adatto al problem solving perchè guidato dall'obiettivo. Esplora solo ciò che è rilevante, quindi il suo costo può essere molto minore di un costo lineare nella dimensione della KB.

\end{document}